\textbf{COLUNI 2025 -- Questão 36}\par\medskip Leia, atentamente, os versos a seguir:

\noindent\begin{minipage}[t]{0.49\linewidth}
Acorda, amor\\
Eu tive um pesadelo agora\\
Sonhei que tinha gente lá fora\\
Batendo no portão, que aflição\\
Era a dura, numa muito escura viatura\\
Minha nossa santa criatura\\
Chame, chame, chame lá\\
Chame, chame o ladrão, chame o ladrão
\end{minipage}\hfill\begin{minipage}[t]{0.49\linewidth}
Acorda, amor\\
Não é mais pesadelo nada\\
Tem gente já no vão de escada\\
Fazendo confusão, que aflição\\
São os homens\\
E eu aqui parado de pijama\\
Eu não gosto de passar vexame\\
Chame, chame, chame\\
Chame o ladrão, chame o ladrão\\
\hfill BUARQUE, Chico. Acorda Amor. In: “\textbf{Sinal Fechado}”, 1974.
\end{minipage}

\par

A canção “Acorda Amor” foi lançada por Chico Buarque em 1974, em um disco intitulado “Sinal Fechado”. Trata-se de uma tentativa bem-sucedida do artista em driblar a censura instituída pela ditadura civil-empresarial-militar (1964-1985). Para enganar os censores, várias músicas do “disco” tiveram autoria de compositores fictícios. No caso de “Acorda Amor”, os compositores seriam Julinho da Adelaide e Leonel Paiva, nomes inventados por Chico Buarque, que estava na mira do regime.

Considerando a letra e a mensagem da canção, bem como seus conhecimentos sobre a ditadura civil-empresarial-militar, cujo golpe completou 60 anos em 2024, é CORRETO afirmar que:
\begin{enumerate}
\item a letra denuncia as violências cometidas pela ditadura e menciona, explicitamente, as ações cometidas pelo Estado durante os “anos de chumbo”.
\item a canção utiliza uma linguagem simbólica e irônica para expressar o clima de medo e perseguição vividos nos chamados “anos de chumbo”.
\item a música exalta os valores morais vigentes no período, defendendo a liberdade de expressão e a segurança nacional, em alta no ano de 1974.
\item a música utiliza uma estrutura poética complexa para criticar a ideologia dominante do período e alegorias não amparadas na realidade.
\end{enumerate}
